\chapter{Introduction}
\label{ch:introduction}

\section{Research Motivation}

Person identification is a recurring requirement in security, surveillance, multimedia search, and human-centred video analytics. A practical investigator may need to answer questions such as whether a person appearing in one video segment is the same person shown in another camera view, whether a still photograph corresponds to a person in a surveillance clip, or whether a natural-language description such as ``a young man wearing a black T-shirt with a white logo'' can retrieve candidate appearances from a video archive. These use cases go beyond conventional closed-set classification because the system must compare visual evidence across different media, resolutions, viewpoints, and temporal contexts.

Traditional person re-identification (Re-ID) research has made substantial progress by learning discriminative visual embeddings for matching images of pedestrians across cameras \citep{zheng2016person,ye2021deep}. However, many Re-ID systems remain difficult to inspect: they return similarity scores but provide limited explanation of which visible traits supported the match. At the same time, modern vision--language models (VLMs) have become capable of producing rich open-domain descriptions of people, clothing, and accessories \citep{radford2021learning,bai2023qwenvl}. This dissertation explores whether VLM-derived attributes can provide a useful bridge between interpretable human descriptions and vector-based retrieval.

\systemname\ is developed as a prototype for this problem. Instead of treating a person crop only as a visual embedding, the system first detects the person, extracts structured visual attributes using a VLM, converts the attributes into both dense semantic vectors and sparse attribute indicators, and then performs identity matching using a weighted hybrid similarity function. The system is designed as an experimental research artifact: it includes a command-line interface, a Streamlit graphical interface, a persistent vector database, and a test dataset creation workflow for extracting person crops from videos.

\section{Problem Statement}

This dissertation studies the following research problem:

\begin{quote}
How can a modular vision--language pipeline represent and match person identities across images and video frames while preserving enough attribute-level structure for inspection, dataset construction, and future text-based retrieval?
\end{quote}

The problem contains three technical tensions. First, visual appearance is high-dimensional and viewpoint-sensitive, so purely symbolic attributes may be too brittle. Second, learned dense embeddings can be effective but often opaque. Third, real-world evaluation requires repeatable datasets and labelled matching pairs, yet early-stage prototypes often lack a convenient mechanism for constructing such datasets from video. The current work responds to these tensions through a hybrid dense--sparse representation and a GUI-supported dataset creation pipeline.

\section{Research Objectives}

The project has five objectives:

\begin{enumerate}
  \item Develop a modular person identification pipeline that processes images and video frames using separate detection, feature extraction, vectorization, and matching modules.
  \item Use a vision--language model to extract structured, interpretable person attributes such as clothing colour, clothing type, accessories, hair style, and distinctive patterns.
  \item Combine dense semantic embeddings with sparse attribute vectors to support both flexible retrieval and inspectable matching evidence.
  \item Implement a persistent vector storage and retrieval layer for incremental identity matching.
  \item Build a GUI-assisted testing workflow that can track persons in video, select trajectories, export reusable test datasets, and run same-person matching checks.
\end{enumerate}

\section{Current Contributions}

At the current stage, the dissertation makes the following contributions.

\begin{itemize}
  \item \textbf{A modular prototype architecture.} \systemname\ separates visual extraction, VLM feature extraction, vectorization, matching, configuration, and interface layers. This makes the system easier to test and extend.
  \item \textbf{A hybrid attribute representation.} The system converts VLM attributes into dense text embeddings and sparse registry-based vectors, then combines cosine similarity and sparse overlap in a configurable matcher.
  \item \textbf{A video-to-testset GUI workflow.} The Streamlit interface supports video upload, person trajectory extraction, manual track selection, crop export, manifest generation, positive/negative pair generation, and ZIP download.
  \item \textbf{A reproducible software artifact.} The repository now contains an Overleaf-ready dissertation draft, a Python package structure, configuration templates without embedded secrets, and lightweight automated tests for dataset export.
\end{itemize}

\section{Dissertation Organization}

This dissertation is organized as a paper-based thesis. \Cref{ch:integrative-methodology} provides the unifying system methodology. \Cref{paper:hybrid-vlm-pid} is the first manuscript and focuses on the core hybrid representation and identity matching architecture. \Cref{paper:video-testset-gui} is the second manuscript and focuses on video-driven test dataset creation and GUI-supported evaluation. \Cref{paper:attribute-registry-protocol} is a third manuscript-in-preparation that formalizes the dynamic attribute registry and outlines the next-stage evaluation protocol. The final chapters synthesize the results, limitations, and future research directions.
